\chapter{2008年上海卷（秋理）}

\section{填空题}

\begin{problem}
不等式 \( |x - 1| < 1 \) 的解集是\fillinblank{}.
\end{problem}

\begin{answer}
\( (0,2) \)
\end{answer}

\begin{problem}
若集合 \(A = \{x \mid x \leq 2\}\)，\(B = \{x \mid x \geqslant a\}\) 满足 \(A \cap B = \{2\}\)，则实数 \(a = \)\fillinblank{}.
\end{problem}

\begin{answer}
\(2\)
\end{answer}

\begin{problem}
若复数 \(z\) 满足 \( z = \i(2 - z) \)（\(\i\) 是虚数单位），则 \(z=\)\fillinblank{}.
\end{problem}

\begin{answer}
\(1+\i\)
\end{answer}

\begin{problem}
若函数 \( f(x) \) 的反函数为 \( f^{-1}(x) = x^2 \)（\(x > 0\)），则 \( f(4) = \)\fillinblank{}.
\end{problem}

\begin{answer}
\(2\)
\end{answer}

\begin{problem}
若向量 \(\bs{a}\)，\(\bs{b}\) 满足 \(|\bs{a}| = 1\)，\(|\bs{b}| = 2\)，且 \(\bs{a}\) 与 \(\bs{b}\) 的夹角为 \(\frac{\pi}{3}\)，则 \(|\bs{a} + \bs{b}| = \)\fillinblank{}.
\end{problem}

\begin{answer}
\(\sqrt{7}\)
\end{answer}

\begin{problem}
函数 \( f(x) = \sqrt{3}\sin x + \sin \left(\frac{\pi}{2} + x\right) \) 的最大值是\fillinblank{}.
\end{problem}

\begin{answer}
\(2\)
\end{answer}

\begin{problem}
在平面直角坐标系中，从六个点： \(A(0,0)\)，\(B(2,0)\)，\(C(1,1)\)，\(D(0,2)\)，\(E(2,2)\)，\(F(3,3)\) 中任取三个，这三点能构成三角形的概率是\fillinblank{}. （结果用分数表示）
\end{problem}

\begin{answer}
\(\frac{3}{4}\)
\end{answer}

\begin{problem}
设函数 \( f(x) \) 是定义在 \( \R \) 上的奇函数. 若当 \( x \in (0, +\infty) \) 时，\( f(x) = \lg x \)，则满足 \( f(x) > 0 \) 的 \( x \) 的取值范围是\fillinblank{}.
\end{problem}

\begin{answer}
\((-1,0)\cup(1,+\infty)\)
\end{answer}

\begin{problem}
已知总体的各个体的值由小到大依次为 \(2, 3, 3, 7, a, b, 12, 13.7, 18.3, 20\)，且总体的中位数为 \(10.5\). 若要使该总体的方差最小，则 \(a\)，\(b\) 的取值分别是\fillinblank{}.
\end{problem}

\begin{answer}
\(a=10.5\)，\(b=10.5\)
\end{answer}

\begin{problem}
某海域内有一孤岛，岛四周的海平面（视为平面）上有一浅水区（含边界），其边界是长轴长为 \(2a\)，短轴长为 \(2b\) 的椭圆. 已知岛上甲，乙导航灯的海拔高度分别为 \(h_1\)，\(h_2\)，且两个导航灯在海平面上的投影恰好落在椭圆的两个焦点上，现有船只经过该海域（船只的大小忽略不计）. 在船上测得甲，乙导航灯的俯角分别为 \(\theta_1\)，\(\theta_2\)，那么船只已进入该浅水区的判别条件是\fillinblank{}.
\end{problem}

\begin{answer}
\(h_1\cot\theta_1+h_2\cot\theta_2\leqslant 2a\)
\end{answer}

\begin{problem}
方程 \(x^{2} + \sqrt{2} x - 1 = 0\) 的解可视为函数 \(y = x + \sqrt{2}\) 的图象与函数 \(y = \frac{1}{x}\) 的图象交点的横坐标. 若 \(x^{4} + ax - 4 = 0\) 的各个实根 \(x_{1}\)，\(x_{2}\)，\(\dots\)，\(x_{k}\)（\(k \leqslant 4\)）所对应的点 \(\left(x_{i}, \frac{4}{x_{i}}\right)\)（\(i=1,2,\dots,k\)）均在直线 \(y = x\) 的同侧，则实数 \(a\) 的取值范围是\fillinblank{}.
\end{problem}

\begin{answer}
\((- \infty,-6)\cup(6,+\infty)\)
\end{answer}

\section{单选题}

\begin{problem}
组合数 \(\mathrm C_n^r\)（\(n > r\geqslant 1\)，\(n,r\in \Z\)）恒等于（\(\quad\)）

\choices
    {\(\frac{r + 1}{n + 1} \mathrm C_{n - 1}^{r - 1}\)}
    {\((n + 1)(r + 1)\mathrm C_{n - 1}^{r - 1}\)}
    {\(nr\mathrm C_{n - 1}^{r - 1}\)}
    {\(\frac{n}{r} \mathrm C_{n - 1}^{r - 1}\)}
\end{problem}

\begin{answer}
D
\end{answer}

\begin{problem}
给定空间中的直线 \(l\) 及平面 \(\alpha\)，条件“直线 \(l\) 与平面 \(\alpha\) 内无数条直线都垂直”是“直线 \(l\) 与平面 \(\alpha\) 垂直”的（\(\quad\)）

\choices
    {充要条件}
    {充分非必要条件}
    {必要非充分条件}
    {既非充分又非必要条件}
\end{problem}

\begin{answer}
C
\end{answer}

\begin{problem}
若数列 \(\{a_{n}\}\) 是首项为\(1\)，公比为 \(a - \frac{3}{2}\) 的无穷等比数列，且 \(\{a_{n}\}\) 的各项和为 \(a\)，则 \(a\) 的值是（\(\quad\)）

\choices
    {\(1\)}
    {\(2\)}
    {\( \frac{1}{2} \)}
    {\( \frac{5}{4} \)}
\end{problem}

\begin{answer}
B
\end{answer}

\begin{problem}
如图，在平面直角坐标系中，\(\Omega\) 是一个与 \(x\) 轴的正半轴，\(y\) 轴的正半轴分别相切于点 \(C\)，\(D\) 的定圆所围成区域（含边界），\(A\)，\(B\)，\(C\)，\(D\) 是该圆的四等分点. 若点 \(P(x, y)\)，点 \(P'(x', y')\) 满足 \(x \leqslant x'\) 且 \(y \geqslant y'\)，则称 \(P\) 优于 \(P'\). 如果 \(\Omega\) 中的点 \(Q\) 满足：不存在 \(\Omega\) 中的其它点优于 \(Q\)，那么所有这样的点 \(Q\) 组成的集合是劣弧（\(\quad\)）

\bitmapfigure[max width=0.62\linewidth]{img/2008/shanghai_science/q15_fig1.png}

\choices
    {\(\myarc{AB}\)}
    {\(\myarc{BC}\)}
    {\(\myarc{CD}\)}
    {\(\myarc{DA}\)}
\end{problem}

\begin{answer}
D
\end{answer}

\section{解答题}

\begin{problem}
如图，在棱长为 \(2\) 的正方体 \(ABCD - A_{1}B_{1}C_{1}D_{1}\) 中，\(E\) 是 \(BC_{1}\) 的中点. 求直线 \(DE\) 与平面 \(ABCD\) 所成角的大小. （结果用反三角函数表示）

\bitmapfigure[max width=0.52\linewidth]{img/2008/shanghai_science/q16_fig1.png}
\end{problem}

\begin{solution}
设 \(F\) 为点 \(E\) 在 \(BC\) 上的射影，连接 \(DF\). 因为 \(E\) 是 \(BC_{1}\) 的中点，所以在正方形 \(BCC_{1}B_{1}\) 中，\(F\) 是 \(BC\) 的中点，从而 \(EF\parallel CC_{1}\). 又因为 \(CC_{1}\perp\text{平面 }ABCD\)，所以 \(EF\perp\text{平面 }ABCD\)，直线 \(DF\) 是直线 \(DE\) 在平面 \(ABCD\) 上的射影，故 \(\angle EDF\) 就是直线 \(DE\) 与平面 \(ABCD\) 所成的角. 由于正方体棱长为 \(2\)，有 \(EF=\dfrac{1}{2}CC_{1}=1\)，\(CF=\dfrac{1}{2}CB=1\)，且 \(DC\perp CF\)，所以 \(DF=\sqrt{DC^{2}+CF^{2}}=\sqrt{5}\). 因而 \(\tan\angle EDF=\dfrac{EF}{DF}=\dfrac{\sqrt{5}}{5}\)，所求角的大小为 \(\arctan\dfrac{\sqrt{5}}{5}\).
\end{solution}

\begin{problem}
如图，某住宅小区的平面图呈圆心角为 \(120^{\circ}\) 的扇形 \(AOB\). 小区的两个出入口设置在点 \(A\) 及点 \(C\) 处，且小区里有一条平行于 \(BO\) 的小路 \(CD\)，已知某人从 \(C\) 沿 \(CD\) 走到 \(D\) 用了 \(10\) 分钟，从 \(D\) 沿 \(DA\) 走到 \(A\) 用了 \(6\) 分钟，若此人步行的速度为每分钟 \(50\) 米，求该扇形的半径 \(OA\) 的长. （精确到 \(1\) 米）

\bitmapfigure[max width=0.5\linewidth]{img/2008/shanghai_science/q17_fig1.png}
\end{problem}

\begin{solution}
由步行速度和用时可得 \(CD=10\times50=500\mathrm{~m}\)，\(DA=6\times50=300\mathrm{~m}\). 因为 \(CD\parallel BO\)，且 \(D,A,O\) 共线，所以 \(\angle CDA=\angle BOA=120^{\circ}\). 在 \(\triangle ACD\) 中，由余弦定理，\(AC^{2}=CD^{2}+AD^{2}-2CD\cdot AD\cos120^{\circ}=500^{2}+300^{2}+2\times500\times300\times\dfrac{1}{2}=700^{2}\)，从而 \(AC=700\mathrm{~m}\). 由余弦定理，\(\cos\angle CAD=\dfrac{AC^{2}+AD^{2}-CD^{2}}{2AC\cdot AD}=\dfrac{11}{14}\). 设圆心 \(O\) 到弦 \(AC\) 的垂线交 \(AC\) 于 \(H\)，则 \(AH=\dfrac{AC}{2}=350\mathrm{~m}\)，且 \(\angle HAO=\angle CAO=\angle CAD\). 在直角三角形 \(AHO\) 中，\(OA=\dfrac{AH}{\cos\angle HAO}=350\times\dfrac{14}{11}=\dfrac{4900}{11}\mathrm{~m}\approx445\mathrm{~m}\). 因此该扇形的半径约为 \(445\mathrm{~m}\).
\end{solution}

\begin{problem}
已知函数 \( f(x) = \sin 2x \)，\( g(x) = \cos \left(2x + \frac{\pi}{6}\right) \). 直线 \(x=t\)（\(t\in\R\)）与函数 \( f(x) \)，\( g(x) \) 的图象分别交于 \(M\)，\(N\) 两点.
\begin{enumerate}
\item 当 \( t = \frac{\pi}{4} \) 时，求 \( |MN| \) 的值；
\item 求 \( \left|MN\right| \) 在 \( t \in \left[0, \frac{\pi}{2}\right] \) 时的最大值.
\end{enumerate}
\end{problem}

\begin{solution}
\begin{enumerate}
\item 当 \(t=\dfrac{\pi}{4}\) 时，\(M\)、\(N\) 的纵坐标分别为 \(f\left(\dfrac{\pi}{4}\right)=\sin\dfrac{\pi}{2}=1\) 和 \(g\left(\dfrac{\pi}{4}\right)=\cos\dfrac{2\pi}{3}=-\dfrac{1}{2}\). 两点的横坐标相同，故 \(\left|MN\right|=\left|1-\left(-\dfrac{1}{2}\right)\right|=\dfrac{3}{2}\).
\item 对任意 \(t\)，由两点横坐标均为 \(t\)，得
\[
\left|MN\right|=\left|\sin2t-\cos\left(2t+\frac{\pi}{6}\right)\right|=\left|\frac{3}{2}\sin2t-\frac{\sqrt{3}}{2}\cos2t\right|=\sqrt{3}\left|\sin\left(2t-\frac{\pi}{6}\right)\right|.
\]
当 \(t\in\left[0,\dfrac{\pi}{2}\right]\) 时，\(2t-\dfrac{\pi}{6}\in\left[-\dfrac{\pi}{6},\dfrac{5\pi}{6}\right]\)，该区间包含使正弦值为 \(1\) 的角 \(\dfrac{\pi}{2}\)，且正弦函数的绝对值不超过 \(1\). 因此 \(\left|MN\right|\) 的最大值为 \(\sqrt{3}\)，在 \(t=\dfrac{\pi}{3}\) 时取得.
\end{enumerate}
\end{solution}

\begin{problem}
已知函数 \( f(x) = 2^x - \frac{1}{2^{|x|}} \).
\begin{enumerate}
\item 若 \( f(x)=2 \)，求 \(x\) 的值；
\item 若 \(2^{t}f(2t) + mf(t) \geqslant 0\) 对于 \(t \in [1,2]\) 恒成立，求实数 \(m\) 的取值范围.
\end{enumerate}
\end{problem}

\begin{solution}
\begin{enumerate}
\item 当 \(x<0\) 时，\(\lvert x\rvert=-x\)，所以 \(f(x)=2^{x}-2^{x}=0\)，不可能等于 \(2\). 当 \(x\geqslant0\) 时，\(f(x)=2^{x}-2^{-x}\)，令 \(u=2^{x}>0\)，则方程化为 \(u-\dfrac{1}{u}=2\)，即 \(u^{2}-2u-1=0\). 因为 \(u>0\)，只能取 \(u=1+\sqrt{2}\)，所以 \(x=\log_{2}\left(1+\sqrt{2}\right)\).
\item 由于 \(t\in[1,2]\)，有 \(f(t)=2^{t}-2^{-t}>0\). 原不等式等价于
\[
2^{t}f(2t)+mf(t)=2^{3t}-2^{-t}+m\left(2^{t}-2^{-t}\right)=\left(2^{t}-2^{-t}\right)\left(2^{2t}+1+m\right)\geqslant0.
\]
由于前一个因式为正，故对每个 \(t\in[1,2]\) 都必须有 \(m\geqslant-(2^{2t}+1)\). 而 \(2^{2t}+1\) 在 \([1,2]\) 上的最小值为 \(5\)，所以恒成立的充要条件是 \(m\geqslant-5\)，即 \(m\in[-5,+\infty)\).
\end{enumerate}
\end{solution}

\begin{problem}
设 \(P(a, b)\)（\(b \neq 0\)）是平面直角坐标系 \(xOy\) 中的点，\(l\) 是经过原点与点 \((1, b)\) 的直线. 记 \(Q\) 是直线 \(l\) 与抛物线 \(x^2 = 2py\)（\(p \neq 0\)）的异于原点的交点. 
\begin{enumerate}
    \item 已知 \(a = 1\)，\(b = 2\)，\(p = 2\)，求点 \(Q\) 的坐标；
    \item 已知点 \(P(a, b)\)（\(ab \neq 0\)）在椭圆 \(\frac{x^2}{4} + y^2 = 1\) 上，\(p = \frac{1}{2ab}\)，求证：点 \(Q\) 落在双曲线 \(4x^2 - 4y^2 = 1\) 上；
\item 已知动点 \(P(a, b)\) 满足 \(ab \neq 0\)，\(p = \frac{1}{2ab}\)，若点 \(Q\) 始终落在一条关于 \(x\) 轴对称的抛物线上，试问动点 \(P\) 的轨迹落在哪种二次曲线上，并说明理由.
\end{enumerate}
\end{problem}

\begin{solution}
\begin{enumerate}
\item 直线 \(l\) 经过原点和点 \((1,b)\)，所以其方程为 \(y=bx\). 在本小题中 \(b=2\)，\(p=2\)，代入抛物线方程得 \(x^{2}=4y=8x\)，即 \(x(x-8)=0\). 除去原点后，\(x=8\)，从而 \(y=2x=16\)，所以 \(Q=(8,16)\).
\item 仍由直线 \(l\) 的方程 \(y=bx\) 和 \(Q\ne(0,0)\)，联立 \(x^{2}=2py\) 得 \(x^{2}=2pbx\)，所以 \(x=2pb\)，且 \(y=2pb^{2}\). 当 \(p=\dfrac{1}{2ab}\) 时，\(Q\) 的坐标为 \(\left(\dfrac{1}{a},\dfrac{b}{a}\right)\). 又因为 \(P(a,b)\) 在椭圆上，所以 \(\dfrac{a^{2}}{4}+b^{2}=1\). 因而
\[
4x^{2}-4y^{2}=4\left(\frac{1}{a}\right)^{2}-4\left(\frac{b}{a}\right)^{2}=\frac{4\left(1-b^{2}\right)}{a^{2}}=\frac{4\cdot a^{2}/4}{a^{2}}=1,
\]
故点 \(Q\) 落在双曲线 \(4x^{2}-4y^{2}=1\) 上.
\item 由上面的联立过程，在本小题中仍有 \(Q=\left(\dfrac{1}{a},\dfrac{b}{a}\right)\). 设这条关于 \(x\) 轴对称的抛物线方程为 \(y^{2}=2q(x-c)\)，其中 \(q\ne0\). 将 \(Q\) 的坐标代入，得
\[
\frac{b^{2}}{a^{2}}=2q\left(\frac{1}{a}-c\right),\qquad b^{2}=2qa-2qca^{2}.
\]
当 \(c=0\) 时，动点 \(P(a,b)\) 的轨迹满足 \(b^{2}=2qa\)，所以轨迹是抛物线.

当 \(c\ne0\) 且 \(qc>0\) 时，将上式整理为
\[
\left(a-\frac{1}{2c}\right)^{2}+\frac{b^{2}}{2qc}=\frac{1}{4c^{2}}.
\]
若 \(qc=\dfrac{1}{2}\)，则上式化为 \(\left(a-\dfrac{1}{2c}\right)^{2}+b^{2}=\dfrac{1}{4c^{2}}\)，轨迹是圆；若 \(qc>0\) 且 \(qc\ne\dfrac{1}{2}\)，轨迹是椭圆.

当 \(qc<0\) 时，上式可写成
\[
\frac{\left(a-\frac{1}{2c}\right)^{2}}{\frac{1}{4c^{2}}}-\frac{b^{2}}{-\frac{q}{2c}}=1,
\]
轨迹是双曲线. 因为题设要求 \(ab\ne0\)，实际轨迹需去掉相应的坐标轴交点，但不影响其二次曲线的类型. 综上，动点 \(P\) 的轨迹可能是抛物线、圆、椭圆或双曲线，具体类型由所给抛物线参数 \(q,c\) 的关系决定.
\end{enumerate}
\end{solution}

\begin{problem}
已知以 \(a_{1}\) 为首项的数列 \(\{a_{n}\}\) 满足： \(a_{n+1} = \begin{cases}
a_{n} + c, & a_{n} < 3,\\
\frac{a_{n}}{d} & a_{n} \geqslant 3.
\end{cases}\)
\begin{enumerate}
\item 当 \(a_{1} = 1\)，\(c = 1\)，\(d = 3\) 时，求数列 \(\{a_{n}\}\) 的通项公式；
\item 当 \(0 < a_{1} < 1\)，\(c = 1\)，\(d = 3\) 时，试用 \(a_{1}\) 表示数列 \(\{a_{n}\}\) 的前\(100\)项的和 \(S_{100}\)；
\item 当 \(0 < a_{1} < \frac{1}{m}\) （\(m\) 是正整数），\(c = \frac{1}{m}\)，正整数 \(d \geqslant 3m\) 时，求证：数列 \(a_{2} - \frac{1}{m}\)，\(a_{3m+2} - \frac{1}{m}\)，\(a_{6m+2} - \frac{1}{m}\)，\(a_{9m+2} - \frac{1}{m}\) 成等比数列当且仅当 \(d = 3m\).
\end{enumerate}
\end{problem}

\begin{solution}
\begin{enumerate}
\item 由递推式可得 \(a_{2}=2\)，\(a_{3}=3\). 因为 \(a_{3}\geqslant3\)，所以 \(a_{4}=\dfrac{a_{3}}{3}=1\)，随后 \(a_{5}=2\)，\(a_{6}=3\). 因此数列每三项重复一次，归纳可得
\[
a_n=\begin{cases}
1,&n=3k-2,\\
2,&n=3k-1,\\
3,&n=3k,
\end{cases}\qquad k\in\Z^+.
\]
\item 当 \(0<a_{1}<1\) 时，前几项为 \(a_{2}=a_{1}+1\)，\(a_{3}=a_{1}+2\)，\(a_{4}=a_{1}+3>3\)，从而 \(a_{5}=\dfrac{a_{1}}{3}+1\). 由此归纳可得，对任意正整数 \(k\)，
\[
a_{3k-1}=\frac{a_{1}}{3^{k-1}}+1,\qquad a_{3k}=\frac{a_{1}}{3^{k-1}}+2,\qquad a_{3k+1}=\frac{a_{1}}{3^{k-1}}+3.
\]
因为每个 \(a_{3k+1}>3\)，除以 \(3\) 后正好得到下一组的第一项. 前 \(100\) 项由 \(a_{1}\) 和 \(33\) 组连续三项组成，故
\[
\begin{aligned}
S_{100}&=a_{1}+\sum_{k=1}^{33}\left(a_{3k-1}+a_{3k}+a_{3k+1}\right)\\
&=a_{1}+\sum_{k=1}^{33}\left(\frac{3a_{1}}{3^{k-1}}+6\right)\\
&=\frac{1}{2}\left(11-\frac{1}{3^{31}}\right)a_{1}+198.
\end{aligned}
\]
\item 先证充分性. 当 \(d=3m\) 时，因 \(a_{1}<\dfrac{1}{m}\)，有 \(a_{3m}=a_{1}+3-\dfrac{1}{m}<3\)，而 \(a_{3m+1}=a_{1}+3>3\)，故
\[
a_{3m+2}=\frac{a_{1}+3}{3m}=\frac{a_{1}}{3m}+\frac{1}{m}.
\]
继续按同一方式计算，得到
\[
a_{6m+2}=\frac{a_{1}}{9m^{2}}+\frac{1}{m},\qquad a_{9m+2}=\frac{a_{1}}{27m^{3}}+\frac{1}{m}.
\]
因此
\[
a_{2}-\frac{1}{m}=a_{1},\quad a_{3m+2}-\frac{1}{m}=\frac{a_{1}}{3m},\quad a_{6m+2}-\frac{1}{m}=\frac{a_{1}}{9m^{2}},\quad a_{9m+2}-\frac{1}{m}=\frac{a_{1}}{27m^{3}},
\]
相邻两项的比均为 \(\dfrac{1}{3m}\)，所以这四项成等比数列.

再证必要性. 当 \(d\geqslant3m+1\) 时，\(a_{3m+1}=a_{1}+3>3\)，且
\[
0<a_{3m+2}=\frac{a_{1}+3}{d}<\frac{1/m+3}{3m+1}=\frac{1}{m}.
\]
于是 \(a_{3m+2}-\dfrac{1}{m}<0\). 从 \(a_{3m+2}<\dfrac{1}{m}\) 出发连续加 \(\dfrac{1}{m}\)，可得 \(a_{6m+2}=a_{3m+2}+3>3\)，从而 \(a_{6m+2}-\dfrac{1}{m}>3-\dfrac{1}{m}>0\). 再除以 \(d\) 后，\(a_{6m+3}=\dfrac{a_{3m+2}+3}{d}<\dfrac{1}{m}\)，同理得到 \(a_{9m+2}=a_{6m+3}+3>3\)，所以 \(a_{9m+2}-\dfrac{1}{m}>0\). 此时第三项与第二项之比为负数，而第四项与第三项之比为正数，不可能成等比数列. 因为 \(d\) 是正整数，\(d\geqslant3m\) 时只有 \(d=3m\) 或 \(d\geqslant3m+1\) 两种情形，故四项成等比数列当且仅当 \(d=3m\).
\end{enumerate}
\end{solution}
